\chapter{現代の欺瞞と主権の回復}
\label{chap:introduction}

\section{設計された認知的不協和}

ソフトウェア工学や物理、数学に携わる技術者が、いわゆる「簿記」や「税務」の教科書を初めて開いたときに覚える激しい違和感と拒絶反応は、決して彼らの知性の欠如に起因するものではない。むしろ逆である。第一原理（First Principles）からの演繹、論理的一貫性、そして保存則の普遍性を信じる訓練を受けた健全な頭脳であるからこそ、そこに蔓延する「設計された不整合」にスタックオーバーフローを起こしているにすぎない。

古くは中世イタリアの商人に起源を持つとされる複式簿記は、本来、代数的な保存則そのものであった。しかし現代の資本主義社会においてそれは、資格制度、士業ギルド、そして高額なサブスクリプションを課すクラウドSaaSベンダーによって幾重もの「認知的トラップ」で包囲され、神聖化されたブラックボックスへと改悪されてしまった。

彼らは、技術者が論理的に世界を理解できないよう、意図的に次のような認知的拷問を仕掛けている。

\begin{enumerate}
    \item \textbf{対称性の意図的な破壊（「借方・貸方」の亡霊）：}\\
    工学や物理において、ノードに対する流入と流出は単なる符号付きスカラー量（流入を $+I$、流出を $-I$）として定式化される。しかし伝統的な簿記教育は、「借方は左、貸方は右」「借りたわけではないが借方」「資産の増加は左だが負債の増加は右」という、言語的・数学的に破綻した呪文の丸暗記を初心者に強要する。これは概念の本質（符号保存則）を隠蔽し、学習者を迷宮に閉じ込めるための第一の障壁である。

    \item \textbf{状態空間の不要な散乱（名詞による目眩まし）：}\\
    10万円未満なら消耗品費、10万円以上30万円未満なら一括償却または少額減価償却資産、30万円以上なら固定資産台帳へ登録して耐用年数テーブルを参照……。これらは数学的には単なる金額と日付を入力とする「閾値判定関数（Step Function）」にすぎない。しかし業界はこれに夥しい名詞のラベルを貼り付け、例外パッチを継ぎ接ぎすることで、「専門家にしか扱えない高等技術」であるかのように演出している。

    \item \textbf{恐怖の煽動と小作農化（Rent-seeking）：}\\
    「少しでも間違えれば税務署に差し押さえられる」「法律は複雑怪奇だから素人が触るな」という恐怖を刷り込むことで、個人の主権を放棄させる。結果として技術者は、銀行口座の平文パスワードを怪しいクラウド家計簿に差し出し、重厚なWeb GUIの前で毎月小作料（月額サブスクリプション）を払い、画面上のボタンをポチポチと押すだけのデジタル小作農へと転落させられる。
\end{enumerate}

\section{第一原理への還元：保存則と指示関数}

この茶番劇の皮を剥ぎ取った後に残る真実は、高校の理数科目レベルの道具立てで完全に記述できるほど極めて質素である。

第一に、あらゆる会計取引の本質は、有向グラフにおけるノード（勘定科目）間のフロー保存則、すなわち\textbf{キルヒホッフの電流則（Kirchhoff's Current Law: KCL）}と同型である。

\begin{equation}
    \sum_{k=1}^n I_k = 0
\end{equation}

勘定科目の集合を頂点、資金の移動を有向枝とする閉鎖系において、出た分だけ入り、全体の代数和は常に厳密にゼロとなる。これ以上でもこれ以下でもない。貸借対照表（B/S）と損益計算書（P/L）の分割すら、この大域的な保存則の移項にすぎない。

\begin{align}
    \text{Assets} + \text{Expenses} &= \text{Liabilities} + \text{Equity} + \text{Revenues} \\
    \sum \text{Assets} + \sum \text{Expenses} - \left( \sum \text{Liabilities} + \sum \text{Equity} + \sum \text{Revenues} \right) &= 0
\end{align}

第二に、世間が「難解な税務」と呼んで恐れおののくものの実態は、取引集合 $T$ に対する\textbf{単項述語（Predicate）に基づく指示関数（Indicator Function）の畳み込み}にすぎない。

\begin{equation}
    \text{Taxable Income} = \sum_{t \in T} \text{Amount}(t) \times \mathbb{I}_{\{\text{is\_taxable}(t) = \text{True}\}} - \text{Deductions}
\end{equation}

国家が課す税法とは、入力された取引データに対して「課税／非課税」「損金算入可／不可」という真偽値（\texttt{true} / \texttt{false}）を返す判定関数にすぎない。この述語論理の判定すら、現代においては最新の税制改正を網羅したフロンティアLLM（Gemini API等）を数円の従量課金で叩く「極小のUNIXフィルタ」として外部委託できる。

\section{主権の奪還}

以上の公理が明らかになった今、我々が取るべき道は自明である。

\begin{itemize}
    \item \textbf{データの不変性：} 財務データはプレーンテキスト（\texttt{.journal}）としてローカルに保持し、Gitの暗号学的コミットで監査ログを自前で確定させる。エンジンには、Haskellの厳密な型システムの上でKCLを保証する \texttt{hledger} を採用する。
    \item \textbf{計算の自立：} 金融シミュレーションや複利計算のためにクラウドを開かない。机の上の物理電卓が備えるわずか2本のハードウェアレジスタ（$M+/M-$ および $GT$）の定数計算、あるいは端末の \texttt{bc -l} によって等比数列を直接解く。
    \item \textbf{パイプラインの純粋性：} 日常の仕訳と税務判定は、\texttt{hledger}、POSIXシェル、\texttt{gawk}、そして単一バイナリのGoフィルタ（\texttt{hledger-ai-taxfilter}）をパイプで繋ぐだけで完結させる。
\end{itemize}

本書は、既存の会計システムが意図的に構築した認知の罠をすべて解体し、技術者が自らの手で経済的・財務的主権を完全掌握するための解毒剤（Antidote）である。

道具の奴隷になるのをやめよ。第一原理に立ち返り、ターミナルから真の自由を宣言しよう。
